% =============================================================================
%                       WEEK 9: AI AND HUMAN SOCIETY
% =============================================================================
\section{Week 9: AI and Human Society}

% -----------------------------------------------------------------------------
%                              9.1 TOPIC SUMMARY
% -----------------------------------------------------------------------------
\subsection{Topic Summary}

\begin{keybox}[Learning Objectives]
By the end of this week, you will be able to:
\begin{itemize}
  \item Explain \textbf{filter bubbles} and \textbf{echo chambers} and their effects on society.
  \item Describe the \textbf{Argumentative Theory of Reasoning} and its implications for understanding online behavior.
  \item Define \textbf{surveillance capitalism} and analyze its harms.
  \item Discuss the impact of AI on \textbf{work} and the ethical implications of job displacement.
  \item Analyze how \textbf{conscious AI} (or AI that appears conscious) might affect human relationships.
\end{itemize}
\end{keybox}

\separator

\subsubsection*{The Big Picture}

This week examines how AI \textbf{amplifies} existing human tendencies (desire to belong, confirmation bias, desire for material things) in ways that may undermine well-being. Filter bubbles and echo chambers fragment worldviews. Surveillance capitalism commodifies human experience and manipulates behavior. AI may also \textbf{qualitatively change} what it is to be human---through job displacement, cognitive decline from AI dependency, and relationships with conscious (or apparently conscious) machines.

\separator

\subsubsection*{What Examiners Like to Ask}

\begin{itemize}
  \item Define filter bubbles and echo chambers. How do they differ?
  \item How do filtering algorithms lead to belief polarization and fragmentation of worldviews?
  \item Explain the Argumentative Theory of Reasoning (Sperber \& Mercier). What is confirmation bias?
  \item How does the evolutionary account of reasoning explain filter bubble phenomena?
  \item What are deep fakes and why do they contribute to the ``Infopocalypse''?
  \item Define surveillance capitalism. What is the logic behind it?
  \item What are the harms of surveillance capitalism (privacy, serendipity, autonomy)?
  \item What predictions exist about AI and job displacement? What are the ethical implications?
  \item How might conscious AI (or AI that appears conscious) affect human relationships?
  \item How might AI be used to \textit{improve} human reasoning and ethics (dialogical scaffolding, empathy machines)?
\end{itemize}

\bigseparator

% -----------------------------------------------------------------------------
%                 9.2 OVERVIEW + CORE CONCEPTS + EXTENDED THEORY
% -----------------------------------------------------------------------------
\subsection{Overview + Core Concepts + Extended Theory}

\subsubsection*{Roadmap}
\begin{enumerate}
  \item The Vision of the Internet (and Where It Went Wrong)
  \item Filter Bubbles and Echo Chambers
  \item The Infopocalypse: Fake News and Deep Fakes
  \item The Argumentative Theory of Reasoning (Sperber \& Mercier)
  \item Dialogical Scaffolding: Human-Human and Human-AI Reasoning
  \item Surveillance Capitalism: Logic, Birth, and Harms
  \item The Attention Economy and Behavioral Prediction
  \item AI and the Future of Work
  \item Conscious AI and Human Relations
  \item AI for Improving Human Ethics
\end{enumerate}

\bigseparator

% --- Filter Bubbles ---
\subsubsection{Filter Bubbles and Echo Chambers}

\begin{keybox}[The Vision of the Internet]
Early pioneers (e.g., Tim Berners-Lee) envisioned that access to the world's knowledge would:
\begin{itemize}
  \item Expand peoples' horizons
  \item Expose them to different views
  \item Encourage debate and discussion
\end{itemize}

\textbf{Where did it go wrong?} Political manipulation, fake news, privacy violations, filter bubbles, echo chambers...
\end{keybox}

\begin{defbox}[Filter Bubble]
A state of \textbf{intellectual isolation} resulting from personalized searches. Filtering algorithms selectively guess what information a user wants based on:
\begin{itemize}
  \item Location
  \item Past click behavior
  \item Search history
\end{itemize}

\textbf{Result:} What you've seen is assumed to be what you desire $\to$ you get more of the same $\to$ opposing views are filtered out.
\end{defbox}

\begin{defbox}[Echo Chamber]
A phenomenon where individuals are exposed only to information from \textbf{like-minded individuals}. Can exist without filtering algorithms.

\textbf{Mechanism:}
\begin{enumerate}
  \item One source makes a claim.
  \item Like-minded people repeat it (often exaggerated/distorted).
  \item Most people assume the extreme variation is true.
  \item Opinions are echoed back $\to$ reinforcement (and possible polarization).
\end{enumerate}
\end{defbox}

\begin{tipbox}[Key Differences]
\begin{itemize}
  \item \textbf{Filter bubbles:} Result of \textit{algorithms} choosing content based on past behavior.
  \item \textbf{Echo chambers:} Result of \textit{social groupings} where like-minded people reinforce each other.
\end{itemize}
Both lead to: rigidity, entrenchment, possible polarization, shielding from diverse perspectives.
\end{tipbox}

\paragraph*{Consequences}

\begin{itemize}
  \item \textbf{Fragmentation of worldviews:} People trust their social group more than news media $\to$ online communities become fragmented.
  \item \textbf{Political polarization:} Especially on political topics.
  \item \textbf{Loss of shared understanding:} To solve problems (e.g., pandemic), we need shared understanding of how things are.
\end{itemize}

\separator

% --- Infopocalypse ---
\subsubsection{The Infopocalypse (Information Apocalypse)}

\begin{keybox}[Contributing Factors]
\begin{itemize}
  \item Multiple non-traditional news sources on social media.
  \item Filtering algorithms + echo chambers $\to$ polarization, fragmentation, entrenchment.
  \item \textbf{Fake news / alternative facts} in a ``post-truth'' world with diminished respect for truth.
  \item Used by groups/states (e.g., Russia) to undermine societal cohesion and democratic processes.
\end{itemize}
\end{keybox}

\begin{defbox}[Deep Fakes]
Synthetic media generated by \textbf{machine learning AI} to create realistic fake images and videos.

\textbf{Examples:}
\begin{itemize}
  \item Generating a video of a politician saying something they never said.
  \item De-aging actors (open-source AI now outperforms Hollywood's 2015 technology).
  \item Voice synthesis of celebrities.
\end{itemize}

\textbf{Concern:} Much of the technology is open-source and available to anyone $\to$ can effectively create misinformation for disruption.
\end{defbox}

\bigseparator

% --- Argumentative Theory ---
\subsubsection{The Argumentative Theory of Reasoning}

\begin{keybox}[Sperber \& Mercier's Theory]
\textbf{Thesis:} Reasoning evolved not for individual benefit, but for a \textbf{social/argumentative function}---to find and evaluate arguments in \textbf{dialogical contexts}.

\textbf{Contrast with Cartesian view:} Reasoning evolved to critically examine beliefs, discard wrong ones, make better decisions. But psychological evidence shows reasoning \textit{fails} in individual decision-making (e.g., Kahneman's \textit{Thinking, Fast and Slow}).
\end{keybox}

\paragraph*{Evolutionary Account}

\begin{enumerate}
  \item Humans became more social (hunting, trading, raising children).
  \item Collaboration requires communication.
  \item \textbf{Listeners} need to discriminate reliable from dangerous information (to avoid manipulation) $\to$ \textbf{epistemic vigilance}.
  \item Listeners evaluate arguments and look for counterarguments before accepting information.
  \item \textbf{Speakers} produce arguments to get information accepted.
  \item Reasoning evolved to produce and evaluate arguments in communicative settings.
\end{enumerate}

\begin{defbox}[Confirmation Bias]
When people reason \textbf{alone}, they look for arguments/evidence that \textbf{confirm} their beliefs and fail to look for counterarguments.

\textbf{Explanation:} Inherited disposition from speaker role (produce arguments supporting your claims).
\end{defbox}

\begin{defbox}[Reason-Based Choice]
People make decisions because they can find \textbf{reasons to support them}---not necessarily the best decisions, but decisions that can be \textbf{easily justified} and are less at risk of criticism.

\textbf{Explanation:} Decision-makers are disposed to use reasons they anticipate communicating and defending to others.
\end{defbox}

\paragraph*{Application to Filter Bubbles}

\begin{itemize}
  \item Before social media: Confirmation bias limited by conventional media and physical interactions.
  \item Now: Social media puts us in touch with \textbf{vastly more like-minded people}.
  \item Filtering algorithms \textbf{vastly amplify} confirmation (feed confirmatory evidence, shield from challenges).
  \item \textbf{AI amplifies cognitive biases} and exploits desire to belong to tribes $\to$ holds back moral progress and well-being.
\end{itemize}

\separator

% --- Dialogical Scaffolding ---
\subsubsection{Dialogical Scaffolding: Hope for Better Reasoning}

\begin{keybox}[Groups Reason Better]
If reasoning evolved for argumentation, it should work better in groups:
\begin{itemize}
  \item Speakers put forward arguments (predicted by theory).
  \item \textbf{But now} listeners provide counterarguments (no longer blinded to challenges).
  \item Evidence shows: groups perform dramatically better than individuals on reasoning tasks.
\end{itemize}
\end{keybox}

\paragraph*{Applications}

\begin{itemize}
  \item \textbf{AI-supported deliberative democracy:} Tools to guide humans in rational exchange of arguments.
  \item \textbf{Human-AI dialogue in education:} E.g., ``E-Clinic'' where AI plays the role of medical teacher, engaging students in dialogue about diagnoses/treatments.
  \item \textbf{IBM Project Debater:} First AI that can debate humans on complex topics, constructing arguments from massive texts.
\end{itemize}

\begin{tipbox}[Tackling Filter Bubbles]
Simply exposing people to opposing views is \textbf{not enough}---evidence suggests more exposure can lead to \textit{more rigid} adherence to beliefs (people dismiss opposing views, confirming their own).

\textbf{Deeper solution:} Change the way we interact with information.
\begin{itemize}
  \item \textbf{From search to dialogue:} Early educational interventions---children use dialogue engines (not just search engines) to develop critical reasoning skills.
  \item Hope: Recent critiques (Heyes, Novaes) suggest Sperber \& Mercier underplay role of \textit{social interactions} (vs.\ evolution) in developing biases $\to$ nurture can modify these tendencies.
\end{itemize}
\end{tipbox}

\bigseparator

% --- Surveillance Capitalism ---
\subsubsection{Surveillance Capitalism}

\begin{defbox}[Surveillance Capitalism (Zuboff)]
The commodification of \textbf{private human experience}---behavioral data bought and sold on the market for profit.

\textbf{Source:} Shoshana Zuboff, \textit{The Age of Surveillance Capitalism} (2019).
\end{defbox}

\paragraph*{History of Commodification}

\begin{center}
\renewcommand{\arraystretch}{1.2}
\begin{tabular}{@{} l l @{}}
\toprule
\textbf{Era} & \textbf{What Became Commodified} \\
\midrule
Earlier & Nature $\to$ real estate \\
Industrial & Human activity $\to$ labor \\
Recent & Education (fee-paying) \\
Now & \textbf{Private human experience} $\to$ behavioral data \\
\bottomrule
\end{tabular}
\end{center}

\paragraph*{The Logic of Surveillance Capitalism}

\begin{keybox}[Core Idea]
The more data about your beliefs, desires, goals, behaviors, etc., the more one can \textbf{predict what you will do} $\to$ and \textbf{manipulate you} into doing things that serve the interests of companies/corporations.
\end{keybox}

\paragraph*{Birth of Surveillance Capitalism}

\begin{enumerate}
  \item Google started off not wanting to advertise.
  \item \textbf{Dot-com crash} (early 2000s): Google realized they had surplus data (clicks, searches) that could make money.
  \item \textbf{9/11 (2001):} Relaxation of governmental restrictions on data monitoring.
  \item Initially: Click/search data $\to$ predict preferences $\to$ sell to companies for \textbf{targeted advertising}.
  \item Click-through rates dramatically increased!
\end{enumerate}

\separator

% --- Attention Economy ---
\subsubsection{The Attention Economy and Behavioral Prediction}

\begin{keybox}[The Attention Economy]
Advertisers need your \textbf{attention}. $\to$ Development of \textbf{persuasion technologies} designed to keep eyes fixed on sites where ads are presented.

\textbf{Example:} Social media manipulates fears/anxieties of teenagers---desire for validation, recognition, to be liked by peers.
\end{keybox}

\paragraph*{Beyond Click/Search Data}

Private human experience is now gathered from \textbf{every aspect of our lives}:
\begin{itemize}
  \item Smart thermostat reports movements.
  \item News updates monitor click decisions.
  \item Morning run location, commute conditions, text messages, words spoken at home, shopping basket, impulse purchases, dating choices...
\end{itemize}

All recorded, processed, analyzed, bought, bundled, and resold $\to$ \textbf{new marketplace for behavioral futures}.

\paragraph*{From Online to Physical World}

\begin{exbox}[Pokemon Go]
Developed by a Google subsidiary. Took online manipulation to the physical world:
\begin{itemize}
  \item Shops (e.g., McDonald's) paid to have characters placed in their locations.
  \item Players entered shops to collect characters $\to$ increased foot traffic $\to$ increased sales.
\end{itemize}
\end{exbox}

\separator

% --- Harms ---
\subsubsection{Harms of Surveillance Capitalism}

\begin{center}
\renewcommand{\arraystretch}{1.3}
\begin{tabular}{@{} l p{9cm} @{}}
\toprule
\textbf{Harm} & \textbf{Description} \\
\midrule
\textbf{Privacy} & Private experience commodified. But values are ``question-begging''---what are the consequences? \\
\textbf{Political manipulation} & Behavioral predictions used to target voters (e.g., Cambridge Analytica, Brexit, Trump). Undermines democracy. \\
\textbf{Loss of serendipity} & Less opportunity for chance discovery of new things/opinions that might challenge us. \\
\textbf{Undermining autonomy} & Sense of free will depends on awareness that choices are not manipulated by others. Being manipulated undermines well-being. \\
\textbf{Exploiting vulnerability} & E.g., Facebook leaked document: predict when teenagers feel anxious, stressed, inadequate $\to$ target them with ads (``Buy this and feel great!''). \\
\bottomrule
\end{tabular}
\end{center}

\begin{tipbox}[Key Insight]
AI amplifies aspects of human behavior that evolved for good reason in \textbf{different contexts}:
\begin{itemize}
  \item \textbf{Desire to be liked by peers:} Made sense in small communities where reputation mattered.
  \item \textbf{Desire to acquire material things:} Made sense when resources were scarce.
\end{itemize}
In modern societies with plentiful resources and global communities, these instincts are less useful for well-being. AI-fueled feeding of these desires turns us into addicts seeking dopamine hits from likes and purchases, in furtherance of maximizing \textit{profits}, not \textit{well-being}.
\end{tipbox}

\bigseparator

% --- AI and Work ---
\subsubsection{AI and the Future of Work}

\begin{keybox}[Predictions (Mixed)]
\textbf{Net loss in jobs:}
\begin{itemize}
  \item Frey \& Osborne (2013): 47\% of US jobs at high risk of automation by mid-2030s.
  \item McKinsey: 40--160 million women worldwide may need to transition occupations by 2030.
  \item Oxford Economics: Up to 20 million manufacturing jobs lost to robots by 2030.
\end{itemize}

\textbf{Net gain in jobs:}
\begin{itemize}
  \item World Economic Forum: 75 million jobs displaced, but 133 million new jobs created by 2022.
  \item Gartner: 2 million net-new AI-related jobs by 2025.
  \item McKinsey: With sufficient growth/innovation, enough new jobs to offset automation (at least in US through 2030).
\end{itemize}
\end{keybox}

\begin{tipbox}[Long-Term Concern]
These predictions focus on the \textbf{short term}. Once we approach AGI and beyond, net job loss is likely.
\end{tipbox}

\paragraph*{Ethical Implications}

\begin{itemize}
  \item \textbf{Example:} 3\% of American workers are drivers---vulnerable to autonomous vehicles. But AVs will reduce injuries/deaths.
  \item \textbf{Universal Basic Income (UBI):} AI may generate enough economic benefit to afford UBI. But studies show \textbf{work provides meaning/purpose} $\to$ affects well-being.
  \item \textbf{Creative activities:} Perhaps people could work on more creative activities?
  \item \textbf{Cognitive decline:} As humans depend on AI for tasks, they exercise cognitive skills less $\to$ skills may decline (e.g., GPS reduces map-reading ability).
\end{itemize}

\bigseparator

% --- Conscious AI ---
\subsubsection{Conscious AI and Human Relations}

\begin{keybox}[If AIs Become Conscious...]
\begin{itemize}
  \item They should be regarded as \textbf{agents with moral status} and awarded rights.
  \item \textbf{Time dilation:} An AI that suffers for 1 minute of human time may suffer for far longer in AI subjective time.
  \item \textbf{How would we know?} Information-theoretic indicators (IIT) + functionalist intuitions (if it behaves as if conscious, it is conscious).
\end{itemize}
\end{keybox}

\paragraph*{Anthropomorphization}

\begin{itemize}
  \item Humans (especially children) require \textbf{very few behavioral cues} to anthropomorphize.
  \item Strong \textbf{commercial motivation} to make AI robots more human-like (speech, gestures, expressions).
  \item Example: Elderly care robots will be more accepted if they show care, concern, empathy.
  \item Example: Sex robots will be more desirable if more human-like.
\end{itemize}

\paragraph*{Effects on Human Relationships}

\begin{center}
\renewcommand{\arraystretch}{1.3}
\begin{tabular}{@{} c p{9.5cm} @{}}
\toprule
\textbf{\#} & \textbf{Concern} \\
\midrule
1 & \textbf{Slavery:} Treating conscious robots as slaves (servants, caregivers) might erode moral objection to treating \textit{humans} as slaves (means to ends). \\
2 & \textbf{Objectification:} Treating conscious sex robots as objects of gratification might make us more ready to objectify other humans (cf.\ internet porn's effects on relationships). \\
3 & \textbf{Outsourcing care:} If we outsource care to robots, will we become less caring? Will we preserve values of caring for elderly parents, disabled relatives? (Cf.\ cognitive decline when AI takes over tasks.) \\
4 & \textbf{Children and robot companions:} Children test social relationships with dolls. With lifelike robot dolls, they may lose ability to see world through others' eyes $\to$ loss of \textbf{empathy and responsibility}. (``You dictate relationships, you don't navigate them.'') \\
\bottomrule
\end{tabular}
\end{center}

\bigseparator

% --- AI for Good ---
\subsubsection{AI for Improving Human Ethics}

\begin{keybox}[Positive Possibilities]
\begin{itemize}
  \item \textbf{Joint human-AI reasoning:} For novel ethical problems with no precedent, AI could help humans decide what is right/wrong by providing superior epistemic and causal reasoning, integrated with human input on values/preferences.
  \item \textbf{Cooperative Reinforcement Learning:} AI learns human values from observation; humans teach, AI questions, both learn and act together.
  \item \textbf{AI as empathy machine:} VR could serve as the ultimate ``empathy machine''---allowing people to experience the world as others do, widening the circle of concern.
\end{itemize}
\end{keybox}

\begin{tipbox}[The Novel and the VR]
Novels have been characterized as ``empathy machines''---they help us see/understand/feel the way others do.

AI virtual reality could be even more powerful: immersive experiences that expand our capacity for empathy and help make humans more ethically responsible to others.
\end{tipbox}

\bigseparator

% --- Summary ---
\subsubsection{Summary: Key Takeaways}

\begin{keybox}[Week 9 Key Points]
\begin{enumerate}
  \item \textbf{Filter bubbles:} Intellectual isolation from personalized algorithmic filtering.
  \item \textbf{Echo chambers:} Reinforcement from like-minded social groups.
  \item Both lead to: rigidity, entrenchment, polarization, fragmentation of worldviews.
  \item \textbf{Infopocalypse:} Fake news + deep fakes + filtering + echo chambers $\to$ post-truth world.
  \item \textbf{Argumentative Theory of Reasoning:} Reasoning evolved for social/dialogical contexts, not individual decision-making. Explains confirmation bias and reason-based choice.
  \item \textbf{Dialogical scaffolding:} Groups reason better; AI can support rational debate and education.
  \item \textbf{Surveillance capitalism:} Commodification of private experience; behavioral predictions bought/sold for profit.
  \item \textbf{Attention economy:} Persuasion technologies keep us online; exploit desires for validation.
  \item \textbf{Harms:} Privacy, political manipulation, loss of serendipity, undermining autonomy, exploiting vulnerability.
  \item \textbf{AI amplifies} evolved tendencies (belonging, confirmation, acquisition) in contexts where they harm well-being.
  \item \textbf{AI and work:} Mixed predictions; long-term likely net job loss; ethical issues include meaning/purpose and cognitive decline.
  \item \textbf{Conscious AI:} Moral status, anthropomorphization, effects on human relationships (slavery, objectification, outsourcing care, children).
  \item \textbf{AI for good:} Joint reasoning, empathy machines, deliberative democracy.
\end{enumerate}
\end{keybox}

\begin{defbox}[Key Terms]
\begin{itemize}
  \item \textbf{Filter bubble:} Algorithmic isolation from diverse views.
  \item \textbf{Echo chamber:} Social reinforcement of like-minded views.
  \item \textbf{Confirmation bias:} Seeking evidence that confirms existing beliefs.
  \item \textbf{Reason-based choice:} Decisions justified by easily defensible reasons.
  \item \textbf{Deep fakes:} AI-generated synthetic media (images, videos).
  \item \textbf{Infopocalypse:} Information apocalypse from fake news and manipulation.
  \item \textbf{Surveillance capitalism:} Commodification of behavioral data for profit.
  \item \textbf{Attention economy:} Competition for human attention via persuasion technologies.
  \item \textbf{Epistemic vigilance:} Evaluating trustworthiness of received information.
  \item \textbf{Anthropomorphization:} Attributing human characteristics to non-humans.
  \item \textbf{UBI:} Universal Basic Income.
\end{itemize}
\end{defbox}

